%% GENERATED by python/paper/build.py from results/ -- do not edit by hand.
%% Rebuild:  .venv\Scripts\python.exe python\paper\build.py --only UK_CRRA_Ageing_vectorX
%% Source:   results/shocks/UK_shocks.csv
\begin{table}[!htb]
\centering
\begin{threeparttable}
\caption{The effect of ageing across the IES, the UK, 2020 (vector $X_i$ calibration)}
\label{table:UK:CRRA:ageing_vectorX}
\renewcommand{\arraystretch}{1.25}
\begin{tabularx}{.9\textwidth}{YYYYY}
\toprule
\textbf{Scenario} & \textbf{CRRA} ($\rho$) & \textbf{Tax rate} & \textbf{Savings rate} & \textbf{Avg. workweek (hours)} \\
\midrule 
 & 0.5 & 11.86\% & 18.92\% & 34.73 \\ % row: baseline@0.5
Baseline & 1.0 & 11.86\% & 17.20\% & 34.73 \\ % row: baseline@1.0
 & 2.0 & 11.86\% & 15.80\% & 34.73 \\[.5em]\hline\\[-.75em] % row: baseline@2.0
 & 0.5 & 14.26\% & $-0.77$ p.p. & 34.66 \\ % row: mild@0.5
Mild ageing & 1.0 & 14.12\% & $-0.77$ p.p. & 34.63 \\ % row: mild@1.0
 & 2.0 & 13.96\% & $-0.77$ p.p. & 34.60 \\[.5em]\hline\\[-.75em] % row: mild@2.0
 & 0.5 & 16.86\% & $-1.57$ p.p. & 34.59 \\ % row: acute@0.5
Acute ageing & 1.0 & 16.59\% & $-1.57$ p.p. & 34.52 \\ % row: acute@1.0
 & 2.0 & 16.19\% & $-1.54$ p.p. & 34.47 \\[.5em]\hline\\[-.75em] % row: acute@2.0
\bottomrule
\end{tabularx}
\begin{tablenotes}[flushleft]
\footnotesize
\item[] \textit{Note:} Every $\rho$ is separately calibrated. Each counterfactual is a separate equilibrium path on which the changed characteristic holds throughout, every other parameter at its value in the UK calibration. The tax rate and the workweek are levels. The savings rate is savings relative to GDP: the baseline rows report its level and every other row the change against the baseline at the same $\rho$, in percentage points. Vector-$X_i$ calibration: see the note to table~\ref{table:US:Calib_vectorX}.
\end{tablenotes}
\end{threeparttable}
\end{table}
